% Este fichero es parte del N�mero 1 de la Revista Occam's Razor
% Revista Occam's Razor N�mero 1
%
% (c) 2007, 2009, Occam's Razor.
%
% Esta obra est� bajo una licencia Reconocimiento 3.0 Espa�a de
% Creative Commons. Para ver una copia de esta licencia, visite
% http://creativecommons.org/licenses/by/3.0/es/ o envie una carta a
% Creative Commons, 171 Second Street, Suite 300, San Francisco,
% California 94105, USA. 

% Seccion Consultorio
%
% Incluye im�genes
\rput(1,-3){\resizebox{!}{4cm}{{\epsfbox{botiquin.eps}}}}
\msection{yellow}{black}{0.2}{CONSULTORIO}

% ---------------------------------------
% Cabecera
\begin{flushright}
\mtitle{6cm}{Preg�ntale a OCCAM}

\msubtitle{13cm}{Todo lo que nunca quiso saber y s� se atrevi� a preguntar}

{\sf por The Occam's Razor Team}

{\psset{linecolor=black,linestyle=dotted}\psline(-10,0)}

\end{flushright}

\vspace{4mm}
% ---------------------------------------

\begin{multicols}{3}

\small
\definecolor{introcolor}{rgb}{0.8,0.8,1.0}
\definecolor{excolor}{rgb}{0.8,0.8,0.8}


\textsf{\textbf{El m�tico N�mero Zero}}
\vspace{2mm}
\hrule
\vspace{2mm}
{\em Sigo vuestra publicaci�n desde el principio y estoy encantada con
los contenidos que inclu�s en cada entrega. Sin embargo, no consigo
encontrar el legendario n�mero Zero de Occam's Razor. He o�do hablar de �l en
varios foros un tanto ominosos, y me encantar�a conseguirlo para
completar mi colecci�n de Occam's Razor.

\begin{flushright}
Una Tauro

Talavera
\end{flushright}
}

\vspace{4mm}
---

Querida Tauro de Talavera, no me seas calavera. El camino hacia el
Zero es tortuoso y lleno de penurias, no apto para
cualquiera. Profundiza en tu interior, y cuando llegues a lo m�s hondo
de tu esencia, entonces, y solo entonces, encontraras el camino hacia
el Zero.

Pero no te entretengas en tu b�squeda. Recuerda que tras el Zero, esta
el menos uno, el menos dos, \ldots

\vspace{6mm}

\textsf{\textbf{Esas criaturas}}
\vspace{2mm}
\hrule
\vspace{2mm}
{\em 
Antes de nada me gustar�a felicitaros por vuestra revista, me est�
resultando muy �til en mi trabajo y espero impaciente cada nuevo
n�mero. Bueno, me gustar�a saber si hay alguna forma para evitar que
se me llene todo de \texttt{zombies}... salen por todas partes... es
una pesadilla... Dios mio!!!!.

\begin{flushright}
VanHelsing

Talansilvania
\end{flushright}
}



\vspace{4mm}
---



Estimado Van, para terminar con tu pesadilla de los no-muertos, puedes
utilizar cualquiera de las dos t�cnicas que describimos a
continuaci�n:

\columnbreak

\begin{enumerate}
\item Ignora las se�ales SIGCHLD utilizando el comando \texttt{signal
(SIGCHLD, SIG\_IGN);}
\item Incorpora un manejador de la se�al SIGCHLD a tu programa para
enterarte de la muerte de cada uno de tus hijos, y as� poder esperar a
que la palmen del todo con un \texttt{wait}
\end{enumerate}



\vspace{6mm}

\textsf{\textbf{Empezando con Linux}}
\vspace{2mm}
\hrule
\vspace{2mm}
{\em 
Querida razor, quiero introducirme en el mundo linux, pero sigo
necesitando usar mi Windows actual, puedo tener las dos cosas en mi
Pc?? 


\begin{flushright}
Jenny

Windowsland
\end{flushright}
}

\vspace{4mm}
---

Estimada Jenny. Claro que puedes tener los dos sistemas en tu PC. Para
ello tienes varias opciones. Quiz�s la mejor para probar si te gusta
es utilizar lo que se llama una versi�n Live-CD con la que podr�s
utilizar linux sin tener que realizar ninguna instalaci�n. Ubuntu o
Knoppix son dos buenas opciones.



Una segunda opci�n es utilizar un emulador de PC como VMware Player o
qemu. La distribuci�n DSL de la que hablamos en este n�mero de la
revista trae qemu integrado, con lo que podr�s ejecutar linux en una
ventana dentro de tu sistema Windows. Sigue las instrucciones que trae
la propia distribuci�n.


\vspace{6mm}

\columnbreak

\textsf{\textbf{Linux sin Linux}}
\vspace{2mm}
\hrule
\vspace{2mm}
{\em 
Una para Occam. Me gustar�a saber si puedo utilizar UNIX y todas
estas cosas tan chulas de las que hablais en la revista sin tener que
instalar uno.


\begin{flushright}
Lina Porgan

Lapataloca
\end{flushright}
}

\vspace{4mm}
---

Hola Lina Porgan. Claro que puedes. Lo m�s sencillo es que instales
Cygwin. Cygwin es un entorno estilo linux para windows que contiene la
mayor�a de las herramientas disponibles en un sistema linux. La
instalaci�n es muy sencilla, solo sigue las instrucciones que te dan
en:

{\tt http://www.cygwin.com/}



\vspace{4cm}

\colorbox{excolor}{
\begin{minipage}{.9\linewidth}
{\bf\sf\Large TIENES ALGUNA\\ DUDA?}
\vspace{1mm}
\hrule
\bigskip

Enviadnos vuestras preguntas tecnol�gicas e intentaremos hacer lo que
podamos, para aclarar cualquier duda.

Pod�is enviarlas a:

\bigskip

{\tt occams-razor@uvigo.es}

\bigskip



\bigskip

\end{minipage}
}

\end{multicols} 


\colorbox{introcolor}{
\begin{minipage}{1.0\linewidth}
\begin{center}
{\bf\sf\Large EVENTOS DE INTER�S}
\vspace{1mm}
\hrule

\includegraphics[height=2.0cm]{fic.eps}

{\Large\bf\sf Jerez de la Frontera. 7, 8 y 9 de Marzo de 2007}

\end{center}
\end{minipage}
}


\vspace{2cm}

\clearpage
\pagebreak
